\documentclass[11pt]{article}
\usepackage[T1]{fontenc}
\usepackage{amsmath,amssymb,amsthm}
\usepackage[margin=1in]{geometry}
\usepackage[hidelinks]{hyperref}
\usepackage{booktabs}

\newtheorem{theorem}{Theorem}
\newtheorem{lemma}[theorem]{Lemma}
\newtheorem{proposition}[theorem]{Proposition}
\newcommand{\Prob}{\mathbb P}

\title{Exact Worst-Case Success Probability for Increasing Event Kernels on
  $K_5$ at $p=1/2$}
\author{Xiaojian Zeng\\
PBC School of Finance, Tsinghua University, 100190\\
\texttt{xiaojianzeng@foxmail.com}}
\date{}

\begin{document}
\maketitle

\begin{abstract}
We determine an exact finite optimization problem arising from a conjecture
about increasing auxiliary events in Bernoulli bond percolation.  On
$K_5$, with one origin and the other four vertices designated as the boundary,
fix $p=1/2$.  If $F(r)$ denotes the least unconditional probability of
connecting the origin to a boundary vertex whose event occurs, among all
jointly defined event families with increasing conditional marginal kernels
and event probabilities at least $r$, then $F$ is an explicitly determined
four-piece affine function on $3/4\leq r\leq1$.  In particular, the bound
$1-2\sqrt\varepsilon$ holds for every $0<\varepsilon\leq1/4$ on this
complete $p=1/2$ slice.  The proof uses lossless symmetry reduction, exact
primal tables at five breakpoints, exact dual certificates on four intervals,
and a verifier that expands the tables over all $1024$ percolation
configurations.  No value $p\neq1/2$ is controlled, so the original all-$p$
problem remains unresolved.
\end{abstract}

\section{Context and result}

Flammant formulated a finite-graph conjecture in which each boundary vertex
is assigned a highly probable event whose conditional probability is
increasing in the percolation configuration \cite{Flammant2026}.  The
conjecture asks for a graph-independent modulus tending to one.  For planar
graphs with the origin and boundary on the outer face, Flammant proved the
explicit modulus $1-2\sqrt\varepsilon$.  That theorem does not apply to
$K_5$.

We study the following finite nonplanar test of that modulus.  Let
\[
 V=\{0,1,2,3,4\},\qquad B=\{1,2,3,4\},\qquad G=K_5,
\]
where $0$ is the origin.  A percolation configuration
$\pi\in\{0,1\}^{E(G)}$ has independent Bernoulli-$p$ coordinates.  Auxiliary
randomness is allowed.  Four jointly defined events $(E_b)_{b\in B}$ are
admissible when
\[
 q_b(\omega):=\Prob(E_b\mid\pi=\omega)
\]
is coordinatewise increasing for every $b$.  No independence is assumed
among the events or between the auxiliary randomness and $\pi$.

For the rest of the paper set $p=1/2$.  For $3/4\leq r\leq1$, let $F(r)$ be
the minimum, over all admissible event families satisfying
$\Prob(E_b)\geq r$ for all $b$, of
\begin{equation}\label{eq:success}
 \Prob\!\left(\bigcup_{b\in B}
   \bigl(\{0\leftrightarrow b\}\cap E_b\bigr)\right).
\end{equation}

\begin{theorem}\label{thm:main}
For $3/4\leq r\leq1$,
\begin{equation}\label{eq:formula}
F(r)=
\begin{cases}
\dfrac{182}{207}r,
  &\dfrac34\leq r\leq\dfrac{207}{256},\\[5pt]
\dfrac{19}{18}r-\dfrac{73}{512},
  &\dfrac{207}{256}\leq r\leq\dfrac{243}{256},\\[5pt]
\dfrac43r-\dfrac{13}{32},
  &\dfrac{243}{256}\leq r\leq\dfrac{63}{64},\\[5pt]
2r-\dfrac{17}{16},
  &\dfrac{63}{64}\leq r\leq1.
\end{cases}
\end{equation}
Consequently, for every $0<t\leq1/2$,
\begin{equation}\label{eq:modulus}
 F(1-t^2)>\frac{15}{16}(1-2t).
\end{equation}
Thus the proposed conditional lower bound $1-2\sqrt\varepsilon$ holds on the
complete $p=1/2$ slice for $0<\varepsilon\leq1/4$.
\end{theorem}

The theorem is a result about one parameter slice.  It neither proves nor
disproves the all-$p$ assertion.  The bounded literature comparison made for
this project located the planar theorem above but no equivalent $K_5$
classification; this is not a claim of priority or of formal open status.

Our symmetry averaging is an instance of fixed-space reduction for symmetric
linear programs \cite{BodiHerrJoswig2013}.  The proof architecture also shares
with recent exact finite-model work the separation of restricted scope,
symmetry-reduced discovery, and exact reconstructive verification
\cite{ZengEtAl2026Finite}.  All details particular to the present problem are
given below.

\section{A lossless finite linear program}

For a configuration $\omega$, let $C_0(\omega)\subseteq B$ be the vertices
connected to $0$.  Record the complete joint conditional law of the auxiliary
events by
\begin{equation}\label{eq:Q}
 Q_\omega(S)=\Prob\bigl(\{b\in B:E_b\text{ occurs}\}=S
                    \mid\pi=\omega\bigr),\qquad S\subseteq B.
\end{equation}
Then
\begin{align}
 Q_\omega(S)&\geq0,&
 \sum_{S\subseteq B}Q_\omega(S)&=1,\label{eq:norm}\\
 q_b(\omega)&=\sum_{S\ni b}Q_\omega(S),&
 q_b(\omega)&\leq q_b(\omega')\quad(\omega\leq\omega').
 \label{eq:mono}
\end{align}
At $p=1/2$, the four probability constraints and the objective are
\begin{align}
 2^{-10}\sum_\omega q_b(\omega)&\geq r \quad(b\in B),
 \label{eq:marginal}\\
 \Phi(Q)&=2^{-10}\sum_\omega\sum_{S\subseteq B}
 Q_\omega(S)\,\mathbf 1\{S\cap C_0(\omega)\neq\varnothing\}.
 \label{eq:objective}
\end{align}

\begin{lemma}\label{lem:lossless}
Minimizing \eqref{eq:objective} subject to
\eqref{eq:norm}--\eqref{eq:marginal} is exactly the optimization defining
$F(r)$.
\end{lemma}

\begin{proof}
Every admissible family gives such a table by \eqref{eq:Q}, and
\eqref{eq:objective} is precisely \eqref{eq:success}.  Conversely, from any
feasible table define a finite probability space on pairs $(\omega,S)$ by
\[
 \Prob((\omega,S))=2^{-10}Q_\omega(S),
\]
let $\pi=\omega$, and declare $E_b$ to occur exactly when $b\in S$.  Its
conditional marginals are the prescribed $q_b$.  Hence every feasible table
is realized by jointly defined auxiliary events; in particular, the
formulation imposes no independence that is absent from the problem.
\end{proof}

It suffices in \eqref{eq:mono} to impose comparisons in which $\omega'$ is
obtained from $\omega$ by opening one edge.  Chaining these cover comparisons
gives every coordinatewise comparison.

\begin{lemma}\label{lem:symmetry}
The minimum is unchanged if $Q$ is required to be invariant under all $24$
permutations of $B$ that fix the origin.
\end{lemma}

\begin{proof}
Let $\Gamma\cong S_4$ act simultaneously on configurations and subsets of
$B$, and average a feasible $Q$ over $\Gamma$.  Normalization,
nonnegativity, marginal monotonicity, and the objective are preserved.  For
each fixed $b$, the unconditional marginal of the averaged table is
\[
 \frac14\sum_{j\in B}\Prob(E_j)\geq r.
\]
Thus averaging maps every feasible table to an invariant feasible table with
the same objective.  The reverse inequality follows because invariant tables
are a subclass of all tables.
\end{proof}

The invariant formulation has one variable for every orbit of a pair
$(\omega,S)$.  Direct orbit construction gives $996$ variables and $90$
configuration orbits.  There are $90$ normalization equations.  After equal
rows are identified, the one-edge marginal comparisons give $984$ distinct
inequalities.  Invariance makes the four unconditional event marginals equal,
so one common marginal constraint suffices.

Write the resulting program in the convention
\begin{equation}\label{eq:LP}
 \min c^{\mathsf T}x,\qquad Ax\leq b(r),\qquad Ex=\mathbf1,
 \qquad x\geq0.
\end{equation}
Only the final entry of $b(r)$ depends on $r$; with the unnormalized marginal
row used here, it is $-1024r$.  A pair $(y,z)$ satisfying
\begin{equation}\label{eq:dualfeas}
 y\leq0,\qquad A^{\mathsf T}y+E^{\mathsf T}z\leq c
\end{equation}
therefore gives the lower bound
\begin{equation}\label{eq:dualbound}
 c^{\mathsf T}x\geq(-1024y_{\rm marg})r+\mathbf1^{\mathsf T}z.
\end{equation}

\section{Exact primal and dual certificate}

The ancillary certificate stores invariant primal tables at five values of
$r$.  Its exactly reconstructed data are as follows.
\begin{center}
\begin{tabular}{ccc}
\toprule
$r$ & objective & nonzero orbit values\\
\midrule
$3/4$       & $91/138$ & $175$\\
$207/256$   & $91/128$ & $90$\\
$243/256$   & $55/64$  & $90$\\
$63/64$     & $29/32$  & $90$\\
$1$         & $15/16$  & $90$\\
\bottomrule
\end{tabular}
\end{center}
It also stores four pairs satisfying \eqref{eq:dualfeas}.  The slopes and
intercepts in \eqref{eq:dualbound} are
\begin{equation}\label{eq:lines}
 \left(\frac{182}{207},0\right),\quad
 \left(\frac{19}{18},-\frac{73}{512}\right),\quad
 \left(\frac43,-\frac{13}{32}\right),\quad
 \left(2,-\frac{17}{16}\right).
\end{equation}

\begin{proposition}\label{prop:certificate}
The five stored primal tables satisfy all constraints exactly, with the
objectives displayed above.  Each stored dual pair satisfies
\eqref{eq:dualfeas} exactly and gives the corresponding affine line in
\eqref{eq:lines}.
\end{proposition}

\begin{proof}[Certificate verification]
The certificate records the edge order, the five sparse primal orbit tables,
and the four sparse dual vectors as rational numbers.  The accompanying
verifier reconstructs the $S_4$ action and all $996$ pair orbits.  For every
endpoint it expands the orbit table over all $1024$ configurations and checks
in rational arithmetic:
\begin{itemize}
\item all $1024$ normalization identities;
\item all four event marginals, each equal to the stated $r$;
\item every one-edge-cover marginal monotonicity comparison; and
\item the objective in \eqref{eq:objective}.
\end{itemize}
For each dual it reconstructs $A$, $E$, and $c$, checks the sign of every
inequality multiplier, and checks all $996$ reduced costs
$c-(A^{\mathsf T}y+E^{\mathsf T}z)$ coefficient by coefficient.  It then
recomputes $-1024y_{\rm marg}$ and $\mathbf1^{\mathsf T}z$.  These checks
return no normalization or monotonicity failure, no positive inequality dual,
and no negative reduced cost.  They also reproduce every rational number in
the two displays above.  Thus the proposition is a finite collection of exact
rational identities and inequalities, independently replayable from the
serialized data.
\end{proof}

We now prove the piecewise formula.  Between any two consecutive endpoint
tables, take their convex interpolation having marginal $r$.  It is feasible,
and its objective is the affine interpolation of the two endpoint objectives.
For that interval the corresponding dual in \eqref{eq:lines} gives exactly the
same affine expression as a lower bound.  The matching upper and lower bounds
prove \eqref{eq:formula} throughout each interval.  The adjacent lines meet at
$207/256$, $243/256$, and $63/64$.  Compactness of the finite feasible
polytope shows that the minimum in the definition of $F$ is attained.

\section{Comparison with the percolation modulus}

Because $B$ contains every vertex other than the origin, the event
$\{0\leftrightarrow B\}$ occurs exactly when at least one of the four edges
incident to $0$ is open.  At $p=1/2$,
\[
 \Prob(0\leftrightarrow B)=1-2^{-4}=\frac{15}{16}.
\]
Put $r=1-t^2$ and define
\[
 D(t)=F(1-t^2)-\frac{15}{16}(1-2t).
\]
On every branch of \eqref{eq:formula}, $D$ is a concave quadratic.  At the
successive branch endpoints
\[
 t=\frac12,\quad\frac7{16},\quad\frac{\sqrt{13}}{16},
 \quad\frac18,\quad0,
\]
its values are
\begin{equation}\label{eq:gaps}
 \frac{91}{138},\quad \frac{19}{32},\quad
 \frac{15\sqrt{13}-10}{128},\quad \frac{13}{64},\quad0.
\end{equation}
The first four are positive; in the radical case this follows from
$225\cdot13>100$.  Concavity therefore proves strict positivity on the first
three closed branch intervals.  On the final branch, which is incident with
$t=0$, direct substitution gives
\[
 D(t)=t\left(\frac{15}{8}-2t\right)>0
 \qquad(0<t\leq1/8).
\]
This proves \eqref{eq:modulus}.  Dividing by $15/16$ yields the desired
conditional lower bound.  Notice that the gap has infimum zero as $t$ tends
to zero, but zero is not attained in the stated domain.

\section{Reproducibility, scope, and assistance}

The decisive ancillary files are
\begin{itemize}
\item \texttt{lp\_p\_half\_parametric\_certificate.json}, containing the five
  primal tables and four duals;
\item \texttt{verify\_lp\_p\_half\_parametric.py}, the exact verifier;
\item \texttt{lp\_p\_half\_parametric.py}, which reconstructs the reduced LP
  and records the discovery procedure; and
\item \texttt{lp\_p\_half\_t\_quarter.py}, which defines the graph, group
  action, orbit routines, and connectivity calculation.
\end{itemize}
Run the verifier with Python 3 from the project root.  Its verification path
does not call an optimizer: all decisive comparisons use
\texttt{fractions.Fraction}.  SciPy/HiGHS was used to discover candidate
vertices and duals, which were rationalized and retained only after exact
checks.  Consequently, floating-point solver output is not part of the proof.

OpenAI Codex assisted with literature search, LP exploration, certificate and
verifier development, adversarial audit organization, and manuscript
preparation.  Language-model output is not mathematical evidence.  The result
rests on the reductions in this paper and on the released exact certificate
checks.  No proof assistant was used.

The restriction $p=1/2$ is essential to the scope of the result.  When $p$
varies, both the objective and event-marginal coefficients acquire nonconstant
polynomial weights in $p$.  The present right-hand-side parametric certificate
does not address those changes.  In particular, it gives no conclusion for
any $p\neq1/2$ and leaves the original all-$p$ problem unresolved.

\bibliographystyle{plain}
\bibliography{references}

\end{document}
